\documentclass[11pt,a4paper]{article}

% ============================================================================
% PACKAGES
% ============================================================================
\usepackage[utf8]{inputenc}
\usepackage[T1]{fontenc}
\usepackage{lmodern}
\usepackage[margin=1in]{geometry}
\usepackage{graphicx}
\usepackage{xcolor}
\usepackage{booktabs}
\usepackage{tabularx}
\usepackage{longtable}
\usepackage{enumitem}
\usepackage{amsmath,amssymb}
\usepackage{listings}
\usepackage{fancyhdr}
\usepackage{titlesec}
\usepackage{tcolorbox}
\usepackage{hyperref}
\usepackage{fontawesome5}

% ============================================================================
% COLORS
% ============================================================================
\definecolor{primaryblue}{RGB}{46,134,171}
\definecolor{darkblue}{RGB}{30,60,90}
\definecolor{accentorange}{RGB}{241,143,1}
\definecolor{successgreen}{RGB}{46,125,50}
\definecolor{dangerred}{RGB}{199,62,29}
\definecolor{codebg}{RGB}{245,245,245}
\definecolor{codeframe}{RGB}{200,200,200}

% ============================================================================
% STYLING
% ============================================================================
\hypersetup{
    colorlinks=true,
    linkcolor=primaryblue,
    urlcolor=primaryblue,
    citecolor=primaryblue,
}

% Headers and footers
\pagestyle{fancy}
\fancyhf{}
\fancyhead[L]{\textcolor{darkblue}{\textbf{GNC Interview Prep Guide}}}
\fancyhead[R]{\textcolor{darkblue}{\thepage}}
\fancyfoot[C]{\textcolor{gray}{\small Miami $\rightarrow$ Moon $\rightarrow$ Jupiter $\rightarrow$ Miami}}
\renewcommand{\headrulewidth}{0.5pt}
\renewcommand{\footrulewidth}{0.3pt}

% Section formatting
\titleformat{\section}
    {\Large\bfseries\color{darkblue}}
    {\thesection}{1em}{}[\titlerule]

\titleformat{\subsection}
    {\large\bfseries\color{primaryblue}}
    {\thesubsection}{1em}{}

\titleformat{\subsubsection}
    {\normalsize\bfseries\color{darkblue}}
    {\thesubsubsection}{1em}{}

% Code listings
\lstset{
    backgroundcolor=\color{codebg},
    basicstyle=\ttfamily\small,
    breaklines=true,
    frame=single,
    rulecolor=\color{codeframe},
    numbers=none,
    showstringspaces=false,
    tabsize=4,
    keywordstyle=\color{primaryblue}\bfseries,
    commentstyle=\color{successgreen}\itshape,
    stringstyle=\color{dangerred},
}

% Custom boxes
\tcbset{
    colback=codebg,
    colframe=primaryblue,
    fonttitle=\bfseries,
    boxrule=0.5pt,
    arc=2pt,
}

\newtcolorbox{keyequation}[1][]{
    colback=blue!5,
    colframe=primaryblue,
    title=Key Equation,
    #1
}

\newtcolorbox{interviewanswer}[1][]{
    colback=green!5,
    colframe=successgreen,
    title=\faComments~Interview Answer,
    #1
}

\newtcolorbox{errorbox}[1][]{
    colback=red!5,
    colframe=dangerred,
    title=\faExclamationTriangle~Error \& Fix,
    #1
}

% ============================================================================
% DOCUMENT
% ============================================================================
\begin{document}

% Title page
\begin{titlepage}
    \centering
    \vspace*{2cm}

    {\Huge\bfseries\color{darkblue} GNC Project\\[0.3cm]
    Interview Preparation Guide}

    \vspace{1cm}

    {\Large\color{primaryblue} Complete Technical Deep-Dive}

    \vspace{2cm}

    \begin{tcolorbox}[colback=blue!5,colframe=primaryblue,width=0.8\textwidth]
        \centering
        \large
        \textbf{Mission:} Miami $\rightarrow$ Moon (2 orbits) $\rightarrow$ Jupiter (3 orbits) $\rightarrow$ Miami\\[0.5em]
        \textbf{Duration:} $\sim$4.1 years\\[0.5em]
        \textbf{Total $\Delta V$:} $\sim$10.5 km/s\\[0.5em]
        \textbf{Spacecraft Mass:} 6,000 kg
    \end{tcolorbox}

    \vspace{2cm}

    {\large Languages \& Tools}\\[0.5em]
    \begin{tabular}{cccc}
        \textbf{Python} & \textbf{C++} & \textbf{MATLAB/Simulink} & \textbf{YAML}\\
        $\sim$15,000 lines & $\sim$3,000 lines & Verification & Configuration
    \end{tabular}

    \vfill

    {\color{gray}\today}
\end{titlepage}

\tableofcontents
\newpage

% ============================================================================
\section{Control Systems}
% ============================================================================

\subsection{PID Controller (Classical)}

\textbf{File:} \texttt{control/attitude\_control.py} -- \texttt{PIDController} class

\textbf{Purpose:} Three-axis attitude control with proportional, integral, and derivative terms.

\begin{keyequation}
\[
\boldsymbol{\tau} = K_p \cdot \mathbf{e} + K_i \cdot \int \mathbf{e} \, dt + K_d \cdot \frac{d\mathbf{e}}{dt}
\]
\end{keyequation}

\textbf{Special Features:}
\begin{itemize}[noitemsep]
    \item \textbf{Anti-windup:} Clamps integrator to prevent saturation
    \item \textbf{Derivative filtering:} Low-pass filter on derivative to reduce noise
    \item \textbf{Per-axis tuning:} Different gains for roll, pitch, yaw
\end{itemize}

\begin{interviewanswer}
``I implemented a PID controller with anti-windup and derivative filtering because raw derivative amplifies high-frequency noise. The anti-windup prevents integrator saturation during large errors.''
\end{interviewanswer}

\subsection{LQR Controller (Optimal)}

\textbf{File:} \texttt{control/attitude\_control.py} -- \texttt{LQRController} class

\textbf{Purpose:} Minimizes a quadratic cost function to find optimal gains.

\begin{keyequation}
\[
\text{Minimize: } J = \int_0^\infty \left( \mathbf{x}^T Q \mathbf{x} + \mathbf{u}^T R \mathbf{u} \right) dt
\]
\[
\text{Feedback law: } \mathbf{u} = -K\mathbf{x}, \quad K = R^{-1}B^T P
\]
\end{keyequation}

\textbf{Implementation Steps:}
\begin{enumerate}[noitemsep]
    \item Build linearized state-space model ($A$, $B$ matrices)
    \item Solve the Continuous Algebraic Riccati Equation (CARE)
    \item Compute optimal gain: $K = R^{-1}B^T P$
\end{enumerate}

\begin{interviewanswer}
``LQR provides optimal full-state feedback. I solve the Riccati equation iteratively since we don't have scipy.linalg in flight software. The $Q$ matrix weights state errors, $R$ weights control effort---I tuned these for 0.1\textdegree~pointing with minimal fuel use.''
\end{interviewanswer}

\subsection{Sliding Mode Control (Robust)}

\textbf{File:} \texttt{control/attitude\_control.py} -- \texttt{SlidingModeController} class

\begin{keyequation}
\[
\mathbf{s} = \boldsymbol{\omega}_{err} + \Lambda \cdot \boldsymbol{\theta}_{err}
\]
\[
\mathbf{u} = \mathbf{u}_{eq} + \mathbf{u}_{sw}, \quad \mathbf{u}_{sw} = -\eta \cdot \text{sat}(\mathbf{s}/\phi)
\]
\end{keyequation}

The \texttt{sat()} function replaces \texttt{sgn()} to reduce chattering.

\begin{interviewanswer}
``Sliding mode is ideal for large-angle slews because it's inherently robust to bounded uncertainties. I use a boundary layer ($\phi$) to prevent chattering---without it, the actuators would oscillate at high frequency.''
\end{interviewanswer}

\subsection{MPC (Model Predictive Control)}

\textbf{File:} \texttt{control/optimal\_control.py} -- \texttt{MPCController} class

\textbf{Key Advantages:}
\begin{itemize}[noitemsep]
    \item Handles constraints (torque limits, rate limits)
    \item Preview of reference trajectory
    \item Receding horizon adapts to changes
\end{itemize}

\begin{interviewanswer}
``MPC is the most flexible controller because it naturally handles constraints. I pre-compute the prediction matrices for efficiency. The receding horizon means we re-plan every step, making it robust to disturbances.''
\end{interviewanswer}

% ============================================================================
\section{Navigation \& State Estimation}
% ============================================================================

\subsection{Extended Kalman Filter (EKF)}

\textbf{File:} \texttt{navigation/ekf.py}

\textbf{State Vector (15 states):}
\[
\mathbf{x} = \begin{bmatrix} \text{pos}(3) & \text{vel}(3) & \text{att\_err}(3) & \text{gyro\_bias}(3) & \text{accel\_bias}(3) \end{bmatrix}^T
\]

\subsubsection{Multiplicative Quaternion Error}

\textbf{Why not put quaternion in state vector?}
\begin{itemize}[noitemsep]
    \item Quaternion has 4 elements but only 3 DOF
    \item Unit-norm constraint violates Gaussian assumption
    \item \textbf{Solution:} Estimate small-angle error (3 states), maintain reference quaternion separately
\end{itemize}

\subsubsection{Joseph Form Covariance Update}

\begin{keyequation}
\[
P = (I - KH)P(I - KH)^T + KRK^T
\]
\end{keyequation}

More numerically stable than $P = (I-KH)P$ for long-duration missions.

\begin{interviewanswer}
``I use the multiplicative quaternion formulation because it respects the unit-norm constraint. After each update, I reset the attitude error by folding it into the reference quaternion. The Joseph form covariance update prevents numerical issues during long-duration missions.''
\end{interviewanswer}

\subsection{Sensors Modeled}

\begin{table}[h]
\centering
\begin{tabular}{lll}
\toprule
\textbf{Sensor} & \textbf{Accuracy} & \textbf{Use Case} \\
\midrule
IMU (gyro) & 0.1 mrad/s bias & Attitude rates (100 Hz) \\
Star Tracker & 5 arcsec & Absolute attitude (10 Hz) \\
Sun Sensor & 0.5\textdegree & Coarse attitude, eclipse detection \\
GPS & 10 m / 0.1 m/s & Near-Earth navigation \\
Deep Space Network & 5 m range & Jupiter cruise navigation \\
\bottomrule
\end{tabular}
\end{table}

% ============================================================================
\section{Orbital Mechanics}
% ============================================================================

\subsection{Key Algorithms}

\subsubsection{Hohmann Transfer}
\begin{keyequation}
\[
\Delta V_1 = \sqrt{\frac{\mu}{r_1}} \left( \sqrt{\frac{2r_2}{r_1+r_2}} - 1 \right)
\]
\[
\Delta V_2 = \sqrt{\frac{\mu}{r_2}} \left( 1 - \sqrt{\frac{2r_1}{r_1+r_2}} \right)
\]
\end{keyequation}

\subsubsection{Bi-elliptic Transfer}
Better than Hohmann when $r_2/r_1 > 11.94$

\begin{interviewanswer}
``I implemented Hohmann for simple orbit raising and Lambert for arbitrary two-point boundary value problems. The Lambert solver is essential for interplanetary trajectory design.''
\end{interviewanswer}

\subsection{RK4 Integrator (C++)}

\textbf{File:} \texttt{dynamics\_engine.cpp}

4th-order accuracy, 4 function evaluations per step:

\begin{lstlisting}[language=C++]
k1 = f(t, y)
k2 = f(t + h/2, y + h*k1/2)
k3 = f(t + h/2, y + h*k2/2)
k4 = f(t + h, y + h*k3)
y_next = y + (h/6)*(k1 + 2*k2 + 2*k3 + k4)
\end{lstlisting}

% ============================================================================
\section{Real-Time Systems (C++)}
% ============================================================================

\subsection{Flight Computer Architecture}

\textbf{File:} \texttt{flight\_computer.h}

\subsubsection{Rate Monotonic Scheduling (RMS)}

Tasks with shorter periods get higher priority. Optimal for fixed-priority periodic tasks.

\textbf{Schedulability bound:} $U \leq n(2^{1/n} - 1)$

\begin{table}[h]
\centering
\begin{tabular}{lcc}
\toprule
\textbf{Task} & \textbf{Rate} & \textbf{Priority} \\
\midrule
Control loop & 100 Hz & Highest \\
Navigation & 10 Hz & Medium \\
Guidance & 1 Hz & Lower \\
Telemetry & 5 Hz & Lowest \\
\bottomrule
\end{tabular}
\end{table}

\subsubsection{Triple Modular Redundancy (TMR)}

Run computation 3 times, take majority vote. Protects against radiation-induced bit flips (SEU).

\begin{interviewanswer}
``I implemented rate-monotonic scheduling because it's provably optimal for periodic tasks. TMR provides radiation tolerance using COTS processors---we run critical calculations 3 times and vote.''
\end{interviewanswer}

\subsection{Memory Pool (Zero-Allocation)}

\textbf{File:} \texttt{memory\_pool.h}

\textbf{Why no malloc in real-time?}
\begin{itemize}[noitemsep]
    \item Non-deterministic timing
    \item Fragmentation
    \item Page faults
\end{itemize}

\begin{lstlisting}[language=C++]
template<typename T, size_t N>
class MemoryPool {
    alignas(64) T blocks_[N];  // Cache-aligned
    T* free_list_;             // Intrusive linked list
};
\end{lstlisting}

\begin{interviewanswer}
``Dynamic memory allocation is forbidden in flight software because malloc has non-deterministic timing. I use pre-allocated memory pools with O(1) allocate/free. The blocks are cache-line aligned for optimal performance.''
\end{interviewanswer}

% ============================================================================
\section{Debugging \& Error Resolution}
% ============================================================================

\begin{errorbox}[title=\faExclamationTriangle~Error 1: PD Controller Sign Error]
\textbf{Symptom:} Attitude diverged---spacecraft spun out of control.

\textbf{Root Cause:} Wrong quaternion error convention. Using $q_{err} = q_{current} \times q_{target}^{-1}$ instead of $q_{err} = q_{target} \times q_{current}^{-1}$.

\textbf{Fix:} Standardized on $q_{err} = q_{target} \times q_{current}^{-1}$ with scalar-positive enforcement.
\end{errorbox}

\begin{errorbox}[title=\faExclamationTriangle~Error 2: Quaternion Normalization Drift]
\textbf{Symptom:} After many integration steps, quaternion norm drifted from 1.0.

\textbf{Root Cause:} Numerical integration accumulates floating-point errors.

\textbf{Fix:} Renormalize after every step + multiplicative EKF formulation.
\end{errorbox}

\begin{errorbox}[title=\faExclamationTriangle~Error 3: Integrator Windup]
\textbf{Symptom:} Controller output saturated and stayed saturated.

\textbf{Root Cause:} Unbounded integrator accumulated huge value.

\textbf{Fix:} Anti-windup clamping: \texttt{integral = clip(integral, -limit, limit)}
\end{errorbox}

\begin{errorbox}[title=\faExclamationTriangle~Error 4: malloc() Timing Jitter]
\textbf{Symptom:} Control loop timing varied from 8ms to 50ms.

\textbf{Root Cause:} Dynamic allocation has non-deterministic timing.

\textbf{Fix:} Pre-allocated memory pools, eliminated all dynamic allocation.
\end{errorbox}

\begin{errorbox}[title=\faExclamationTriangle~Error 5: EKF Numerical Instability]
\textbf{Symptom:} Covariance matrix became non-positive-definite after days.

\textbf{Root Cause:} Standard update $P = (I-KH)P$ is numerically unstable.

\textbf{Fix:} Joseph form: $P = (I-KH)P(I-KH)^T + KRK^T$
\end{errorbox}

\begin{errorbox}[title=\faExclamationTriangle~Error 6: Propellant Budget Exceeded]
\textbf{Symptom:} Mission requires $\sim$10.5 km/s but only $\sim$2 km/s available.

\textbf{Root Cause:} Direct transfer to Jupiter isn't mass-feasible with chemical propulsion.

\textbf{Fix:} Proposed VEEGA gravity assists, documented constraint.
\end{errorbox}

% ============================================================================
\section{Quick Reference Card}
% ============================================================================

\begin{tcolorbox}[colback=blue!5,colframe=primaryblue,title=Equations to Know]
\begin{align*}
\textbf{Rocket Equation:} \quad & \Delta V = I_{sp} \cdot g_0 \cdot \ln(m_0/m_f) \\[0.5em]
\textbf{Vis-Viva:} \quad & v^2 = \mu\left(\frac{2}{r} - \frac{1}{a}\right) \\[0.5em]
\textbf{Quaternion Kinematics:} \quad & \dot{q} = \frac{1}{2} q \otimes \omega_q \\[0.5em]
\textbf{Euler's Equations:} \quad & I\dot{\omega} = \tau - \omega \times (I\omega) \\[0.5em]
\textbf{Kalman Gain:} \quad & K = PH^T(HPH^T + R)^{-1}
\end{align*}
\end{tcolorbox}

\vspace{1cm}

\begin{center}
\begin{tcolorbox}[colback=green!5,colframe=successgreen,width=0.8\textwidth]
\centering
\Large\bfseries
Good Luck Tomorrow!\\[0.5em]
\normalsize\normalfont
Remember: You built this. You understand it. You debugged it.\\
\textbf{Be confident.}
\end{tcolorbox}
\end{center}

\end{document}
